% !TeX root = ../main.tex
\section{Fazit und Ausblick}
\label{sec:schluss}

\subsection{Beantwortung der Fragestellungen}
\textbf{Nationale Prägung und Gruppen.} Das Netz ist stark national geprägt.
43,1\,\% der Kanten im Arbeitsnetz verbinden Landsleute, bei zufällig
verbundenen Athleten wären es 1,2\,\%. Louvain findet eine ausgeprägte
Community-Struktur mit zwei Typen. Elf grosse Communities bestehen fast nur
aus einer Nation. Vier grosse internationale Communities trennen sich nach
Generationen und bilden sich um die Meisterschaften ihrer Zeit, von den
Europameisterschaften der frühen 2000er-Jahre bis zu den aktuellen
Weltmeisterschaften.

\textbf{Zentrale Athleten und Teams.} Ein hoher Degree entsteht bei uns
durch wenige, sehr grosse internationale Veranstaltungen und nicht durch
viele Wettkämpfe. Hohe Betweenness kennzeichnet Athleten an den Übergängen
zwischen Generationen oder zwischen nationaler und internationaler Ebene.
Stärkere Schwimmer sind signifikant zentraler, der Zusammenhang ist mit
$\rho$ zwischen 0,17 und 0,36 aber schwach bis mässig. Nationalteams, deren
Athleten sich vor allem mit Landsleuten treffen, erreichen weniger Athleten
anderer Nationen (Spearman $-0{,}82$).

\textbf{Zusammenhalt.} Wiederholte Begegnungen sind kein Zufall. Das
Arbeitsnetz enthält rund zehnmal so viele wiederholte Begegnungen wie das
bipartite Configuration Model. Zusammengehalten wird es von internationalen
Meisterschaften. Fehlen die internationalsten Veranstaltungstage mit 30\,\%
der Teilnahmen, bleiben nur 53,7\,\% der Athleten in der Giant Component,
bei zufälligen Tagen 91,1\,\%. Das Entfernen der zentralsten Athleten
verändert das Netz dagegen kaum.

\textbf{Bevorzugte Verbindungen.} Am stärksten verbinden sich Gleichaltrige
($r = 0{,}85$), danach Landsleute ($r = 0{,}42$). Starke Schwimmer treffen
sich vor allem international untereinander, in der jüngsten Generation am
deutlichsten. National treffen sich dagegen Gleichaltrige unabhängig von
ihrem Niveau.

\textbf{Vorhersage.} Netzwerkverfahren sagen neue Begegnungen fünf- bis
sechsmal besser voraus als der Zufall. Die einfache Regel gleiche Nation ist
aber noch besser. Unter den aktiven Athleten trifft sich mehr als die
Hälfte der Landsleute, die sich im Trainingszeitraum nicht begegnet sind,
im Testzeitraum.

Zusammen ergibt sich ein Netz mit zwei Ebenen. National entsteht ein dichtes,
wiederkehrendes Netz unter Gleichaltrigen. International verbinden wenige
grosse Meisterschaften die Nationen und Generationen, und dort treffen sich
die stärksten Schwimmer.

\subsection{Grenzen, Ausblick und Reflexion}
Die Ergebnisse beruhen auf den persönlichen Bestleistungen der TOP100-Listen
mit Stand 2025. Pro Disziplin ist nur die schnellste Zeit der Karriere gespeichert. Verbessert ein Athlet sie, verschwindet der frühere Tag aus den Daten. Ältere Jahre sind auch deshalb dünner erfasst, besonders vor 2005. Eine Kante zeigt einen gemeinsamen Bestleistungstag und keine persönliche Beziehung, und viele tatsächliche Begegnungen an Wettkämpfen fehlen. Zudem betreffen die 126 fehlenden Profile fast
nur Frauen, und die Profilnation entspricht dem Stand 2025.

Für eine Weiterführung sehen wir mehrere Ansätze.
\begin{itemize}
  \item \textbf{Vollständige Ergebnislisten} statt Bestleistungen würden ein
  echtes Begegnungsnetz liefern. Eine Kante könnte dann auch ein gemeinsames
  Rennen oder Finale bedeuten.
  \item \textbf{Zeitliche Netze pro Saison} würden zeigen, wie sich
  Generationen bilden, wie Karrieren verlaufen und wie stark die Pandemie
  2020 das Netz verändert hat.
  \item \textbf{Netzwerk und Attribute kombiniert} könnten die Link
  Prediction verbessern, etwa Soundarajan-Hopcroft mit Nationen statt
  Louvain-Communities.
  \item \textbf{Vereine als weitere Relation} würden Trainingsgruppen
  sichtbar machen, sobald die Vereinsangaben vollständiger sind.
  \item \textbf{Diffusion} könnte untersuchen, wie sich Neuerungen wie die
  Hightech-Anzüge von 2008 und 2009 über gemeinsame Wettkämpfe ausbreiten.
  \textbf{Graphlets} könnten die Struktur nationaler Netze vergleichen.
\end{itemize}

Gut funktioniert hat, die Analysen aufeinander aufzubauen. Befunde wie der
Kern um die Weltmeisterschaft 2009 oder die nationale Prägung tauchten in
mehreren Kapiteln wieder auf und liessen sich so gegenseitig prüfen.

Anders machen würden wir vor allem die Datengrundlage. Wir haben einen
bestehenden Datensatz aus einem früheren Projekt übernommen, der pro
Disziplin nur die schnellste Zeit der Karriere enthält. Das Netz bildet
deshalb gemeinsame Bestleistungstage ab und nicht alle Begegnungen an
Wettkämpfen. Fehlende Profile liessen sich zudem nicht mehr nachladen, da
Cloudflare den Zugriff blockierte. Mit einer eigenen, auf die Fragestellung
ausgerichteten Erhebung würden wir vollständige Ergebnislisten der
Wettkämpfe sammeln.

Methodisch haben wir vor allem aus der Projektion gelernt. Jeder
Veranstaltungstag erzeugt eine Clique, weshalb Masse wie Clustering, k-Core
oder ein Vergleich mit dem gewöhnlichen Configuration Model einen passenden
Bezugspunkt brauchen. Das bipartite Configuration Model würden wir früher
einsetzen, auch für den Modularity-Test in Kapitel~\ref{sec:communities}.
Kapitel~\ref{sec:assortativitaet} hat zudem gezeigt, wie leicht
zusammengefasste Gruppen einen Effekt vortäuschen. Ergebnisse auf
Teilgruppen zu prüfen, würden wir von Anfang an einplanen.

